\papersection{Backtest}{sec:backtest}

\subsection{What did we test}\label{sec:whattest}

We trade the model of the Model part, with its setting and rule fixed in advance, and score it on
the held-out years. Returns are in bps of one contract's notional, summed rather than compounded,
on a daily grid of trading sessions with zero on days without a trade; return, volatility and
the Sharpe ratio are annualized over 252 sessions. Costs are charged per side at a quarter or a
half tick.

\subsubsection{Time period}

\paragraph{Train, validation and test split.} Table~\ref{tab:split} shows how the years are used.
The first three years only fill the library of past reactions and the history the universe
screen needs. Settings are chosen on the training years, checked on the validation years, and
the test years are used to report results. The gate cut-offs and the size reference come from
2013--2020, so nothing in the test years is fitted.

\begin{table}[!tb]
  \centering\footnotesize
  \caption{Sample split.}\label{tab:split}
  \begin{tabular}{lll}
    \toprule
    Period & Years & Use \\
    \midrule
    Burn-in    & 2010--2012 & library of past reactions and screen history only \\
    Train      & 2013--2017 & choosing the DTW setting, gate and sizing \\
    Validation & 2018--2020 & checking that the choice holds \\
    Test       & 2021--2026 & results \\
    \bottomrule
  \end{tabular}
\end{table}

\paragraph{Why 30 minutes.} Each trade is held for 30 minutes and each path is the 31 minutes
ending at the entry. The horizon matches how long the news keeps moving prices: after hot CPI
prints \ES\ falls about 37 bps in the first minute and a further 17 bps by $T+30$, and is roughly
flat over the following half hour; \NQ\ and \YM\ behave the same way. A 30-minute path is long
enough to contain the whole first reaction and short enough that the library of past releases at
the same clock time holds hundreds of comparable paths. Equal path and holding lengths also make
the five windows of the ladder back-to-back, so no minute is traded twice.

\subsubsection{Assets}

Each asset is first traded on its own with the same setting, gate and sizing: \ES\ and \NQ, the
core of the book; \YM, run through the identical pipeline after the \ES\ and \NQ\ results were
known, as an out-of-sample market; and \ZN. Two multi-asset books follow.

\subsubsection{Equities book}

The equities book combines \ES\ and \NQ\ with the inverse-volatility weights and cost gate of
equation~\eqref{eq:assetweights}; both always pass the gate at a quarter tick, and the weights
average 58\% \ES\ and 42\% \NQ. We also report the three-index book \ES\ + \NQ\ + \YM, with
weights of about 36\%, 26\% and 38\%.

\subsubsection{Combined book}

The combined book adds \ZN\ on the same terms, and the cost gate decides whether it is traded.
Before costs every asset passes and \ZN\ takes about 60\% of the weight; at a quarter tick a round
trip is 10--17\% of \ZN's typical 30-minute move, so it is excluded except in part of 2023
(\ref{sec:treasury_result}).

\subsection{General Statistics}\label{sec:signalquality}

\paragraph{Signal quality.} Before trading, we ask whether the two features rank what happened
next (Figure~\ref{fig:ic}). The direction forecast $\mu$ ranks the next return
\emph{negatively} in the training years ($-0.060$ in \ES, $-0.042$ in \NQ) and \emph{positively}
in the validation years ($+0.069$ and $+0.091$); in \YM\ and \ZN\ it is close to zero in both
($-0.006$ and $+0.018$; $-0.003$ and $-0.005$). The dispersion $\sigma$, in contrast, ranks the
absolute size of the next move positively in every market and period, at 0.10--0.24. The
analogues know how large the next move will be; what they say about its direction is small and
changed sign around 2018.

\begin{figure}[!tb]
  \centering
  \includegraphics[width=0.55\linewidth]{dtw_ic}
  \caption{Information coefficients of the DTW features on admitted games, \ES\ and \NQ. Panel (a):
  rank correlation of $\mu$ with the next 30-minute return. Panel (b): rank correlation of $\sigma$
  with the absolute next 30-minute return.}
  \label{fig:ic}
\end{figure}

\paragraph{Results.} Table~\ref{tab:results} reports the test years. The \ES\ book trades 1,202
times, about 225 a year, and earns a Sharpe ratio of 0.78 before costs, 0.64 at a quarter tick and
0.49 at a half tick, against 0.69 for holding \ES, with a drawdown of $-5.0\%$ against $-31.1\%$
and a beta of 0.007. Half its trades win and 53\% are long, so it is not a disguised long
position. \NQ\ is positive but weak (0.11 at a quarter tick). \YM\ does not work: only five
families pass its screen, and it loses before costs ($-0.32$) and after them ($-0.39$). The
equities book earns 0.54 with a drawdown of $-5.3\%$; adding \YM\ lowers it to 0.27. The combined
book earns 0.40 (Figure~\ref{fig:sharpe}).

\begin{table}[!tb]
  \centering\footnotesize
  \caption{General statistics of the DTW event book, test years 2021--2026, quarter tick per side.}\label{tab:results}
  \begin{threeparttable}
  \setlength{\tabcolsep}{4pt}
  \begin{tabular}{lccccccccc}
    \toprule
    & & \multicolumn{3}{c}{Sharpe ratio} & Buy-and- & Total & Ann.\ & Max & Hit \\
    \cmidrule(lr){3-5}
    Book & Trades & Gross & $\tfrac14$ tick & $\tfrac12$ tick & hold & return & vol & DD & rate \\
    \midrule
    \ES & 1,202 & 0.78 & 0.64 & 0.49 & 0.69 & $10.9\%$ & 3.2\% & $-5.0\%$ & 50.0\% \\
    \NQ & 1,141 & 0.14 & 0.11 & 0.08 & 0.58 & $2.1\%$ & 3.6\% & $-8.2\%$ & 49.5\% \\
    \YM & 870 & $-0.32$ & $-0.39$ & $-0.47$ & 0.67 & $-4.8\%$ & 2.4\% & $-8.0\%$ & 50.0\% \\
    \ZN & 1,945 & 0.34 & $-0.97$ & $-2.24$ & $-0.69$ & $-7.4\%$ & 1.4\% & $-8.0\%$ & 46.9\% \\
    Equities (\ES\ + \NQ) & 2,343 & 0.66 & 0.54 & 0.42 & 0.65 & $7.4\%$ & 2.6\% & $-5.3\%$ & 49.8\% \\
    \ES\ + \NQ\ + \YM & 3,213 & 0.39 & 0.27 & 0.14 & 0.68 & $2.8\%$ & 2.0\% & $-3.0\%$ & 49.8\% \\
    Combined (\ES\ + \NQ\ + \ZN) & 2,476 & 0.76 & 0.40 & 0.41 & 0.24 & $5.4\%$ & 2.5\% & $-5.3\%$ & 49.3\% \\
    \bottomrule
  \end{tabular}
  \begin{tablenotes}\footnotesize
    \item Notes: Direction $\operatorname{sign}(\mu)$, outer 60\% of $\mu$ with cut-offs from
    2013--2020, size $\min(\tilde\sigma/\sigma,2)$, five windows per admitted release. Return,
    volatility, drawdown and hit rate at a quarter tick. Multi-asset books use the weights and
    cost gate of equation~\eqref{eq:assetweights}, so each cost column is the book a trader facing
    that cost would run. Buy-and-hold drawdowns: \ES\ $-31.1\%$, \NQ\ $-48.3\%$, \YM\ $-25.6\%$.
  \end{tablenotes}
  \end{threeparttable}
\end{table}

\begin{figure}[!tb]
  \centering
  \begin{subfigure}[b]{0.58\linewidth}\centering
    \includegraphics[width=\linewidth]{sharpe_by_cost_with_ym}
    \caption{Sharpe ratio by cost level, against buy-and-hold.}\label{fig:sharpe}
  \end{subfigure}\hfill
  \begin{subfigure}[b]{0.4\linewidth}\centering
    \includegraphics[width=\linewidth]{robustness_es}
    \caption{\ES\ Sharpe ratio at a quarter tick by $K$, gate width and cut-off rule.}\label{fig:robust}
  \end{subfigure}
  \caption{The DTW event book, 2021--2026. Panel (b): cut-offs are either fixed from 2013--2020
  or recomputed each year.}
  \label{fig:book}
\end{figure}

\paragraph{Checks.}\label{sec:robust} Table~\ref{tab:checks} reports five checks. Rebuilding the
features with the path running five minutes \emph{past} the entry raises the Sharpe ratio to
about 2.8 in every market, so the real book's timing is sound. Drawing the direction of each trade
at random beats the \ES\ book only 4.1\% of the time, but beats \NQ\ and \YM\ often ($p=0.37$ and
0.76). Holding the same windows always long loses in every book, so the \ES\ profit comes from
direction. \ES\ breaks even at 1.36 ticks per side. Across twelve gate, cut-off and sizing rules
\ES\ and the equities book are always positive, \YM\ only once, and across the number of
neighbours $K$ (Figure~\ref{fig:robust}) the \ES\ Sharpe ratio ranges from 0.17 to 0.64. Broader
searches over about 138,000 settings confirm that the direction is unstable across periods (the
Sharpe ratios of the same settings on 2014--2017 and 2018--2020 correlate at 0.15). The appraisal
ratio---alpha on buy-and-hold over residual volatility---is 0.61 for \ES\ and 0.53 for the
equities book; with a directional IC of 0.027 on about 225 bets a year, the fundamental law of
active management \citep{Grinold2000} implies about 0.4, so 0.4--0.6 is a fair range.

\begin{table}[!tb]
  \centering\footnotesize
  \caption{Checks on the DTW event book, 2021--2026, quarter tick per side.}\label{tab:checks}
  \begin{threeparttable}
  \begin{tabular}{lccccc}
    \toprule
    & \ES & \NQ & \YM & \ES\ + \NQ & \ES\ + \NQ\ + \YM \\
    \midrule
    Leak test: Sharpe with the path 5 minutes past entry & 2.86 & 2.85 & 2.84 & 3.62 & 3.83 \\
    Random-sign test: $p$-value, 5,000 draws & 0.041 & 0.37 & 0.76 & 0.064 & 0.18 \\
    Always long in the same windows: total return & $-2.4\%$ & $-1.0\%$ & $-6.2\%$ & $-1.8\%$ & $-3.6\%$ \\
    Break-even cost, ticks per side & 1.36 & 1.10 & --- & 1.35 & 0.76 \\
    Rules with Sharpe $>0$, of 12 & 12 & 7 & 1 & 12 & 12 \\
    Appraisal ratio vs buy-and-hold & 0.61 & 0.11 & $-0.38$ & 0.53 & 0.26 \\
    \bottomrule
  \end{tabular}
  \begin{tablenotes}\footnotesize
    \item Notes: \YM\ loses before costs, so it has no break-even cost. The \YM\ columns come from
    the same code run with \YM\ added, in which the \ES, \NQ\ and equities results reproduce exactly.
  \end{tablenotes}
  \end{threeparttable}
\end{table}

\subsection{Hit rate per window}\label{sec:hitrate}

Table~\ref{tab:window} and Figure~\ref{fig:byanchor} split the book by window. The first window,
entering at $T+1$, is the best in every equity index: it wins 53\% of its trades in \ES, 53\% in
\NQ\ and 58\% in \YM, with Sharpe ratios of 0.95, 0.38 and 0.41, against 1.7\%, 5.4\% and $-3.0\%$ for
always long in the same windows. The hit rate then falls toward a coin flip in the second to
fourth windows, and the fifth, at $T+121$, is the worst everywhere: 43\% in \ES\ (Sharpe $-0.91$),
46\% in \NQ\ and 45\% in \YM. In \YM\ the first window is the only one that makes money, which is
why its five-window book loses; in \ZN\ no window reaches 50\%. The path carries information while
it still contains the reaction to the release, and less once the reaction is over. Choosing the
window on these test years would be choosing on the test period, so the book keeps all five.

\begin{table}[!tb]
  \centering\footnotesize
  \caption{Hit rate and Sharpe ratio by window, 2021--2026, quarter tick per side.}\label{tab:window}
  \begin{threeparttable}
  \begin{tabular}{lccccc}
    \toprule
    Entry & $T+1$ & $T+31$ & $T+61$ & $T+91$ & $T+121$ \\
    \midrule
    \ES\ hit rate (Sharpe) & 53.4\% (0.95) & 50.2\% (0.27) & 52.1\% (0.33) & 50.0\% (0.45) & 43.4\% ($-0.91$) \\
    \NQ\ hit rate (Sharpe) & 53.3\% (0.38) & 49.3\% (0.02) & 50.0\% ($-0.18$) & 48.5\% ($-0.04$) & 45.8\% (0.01) \\
    \YM\ hit rate (Sharpe) & 58.1\% (0.41) & 47.8\% ($-0.40$) & 46.2\% ($-0.17$) & 52.1\% ($-0.18$) & 44.7\% ($-0.57$) \\
    \ZN\ hit rate (Sharpe) & 47.3\% ($-0.39$) & 48.5\% ($-0.41$) & 46.9\% (0.02) & 45.8\% ($-0.64$) & 45.7\% ($-0.94$) \\
    \bottomrule
  \end{tabular}
  \begin{tablenotes}\footnotesize
    \item Notes: Trades per window: 150--262 in the equity indices and 339--442 in \ZN.
  \end{tablenotes}
  \end{threeparttable}
\end{table}

\begin{figure}[!htb]
  \centering
  \includegraphics[width=0.72\linewidth]{pnl_curves_by_anchor_with_ym}
  \caption{Cumulative net P\&L at a quarter tick by window, 2021--2026, against always long in the
  same windows. Rows: \ES, \NQ, \YM, \ZN; columns: entry at $T+1$, $T+31$, $T+61$, $T+91$ and $T+121$.}
  \label{fig:byanchor}
\end{figure}

\subsection{2022 for equities}\label{sec:2022}

Table~\ref{tab:byyear} splits the test period by year. The \ES\ and equities books are positive
in five of six years; 2022, the year of the inflation shock and the fastest monetary tightening
in four decades, is the exception (equities $-0.87$, $-2.5\%$). A month-by-month breakdown shows
that 2022 was not a period of poor matches: the mean weighted similarity of the neighbours was
about 0.51 in \ES\ and 0.50 in \NQ\ in the first half of 2022, the same as in other years. In those
months the release windows themselves fell (about $-10$ bps on average in \ES\ and $-12$ bps in
\NQ), the \ES\ book was long in about 63\% of its trades, and the worst month, May, was a short
position into a rebound. The paths matched as closely as usual; what those paths had done next in
earlier years no longer held, because the market's reading of the same data had changed with the
Federal Reserve's reaction function. \YM\ moved the other way: it gained in 2022 (0.34) and lost
most in 2025. Win rates are 47--54\% in every year, so the edge, where it exists, comes from
winners being larger rather than more frequent.

\begin{table}[!tb]
  \centering\footnotesize
  \caption{Sharpe ratio by year, quarter tick per side.}\label{tab:byyear}
  \begin{threeparttable}
  \begin{tabular}{lcccccc}
    \toprule
    & 2021 & 2022 & 2023 & 2024 & 2025 & 2026 \\
    \midrule
    \ES & 0.67 & $-0.48$ & 0.89 & 1.33 & 0.76 & 1.27 \\
    \NQ & 0.12 & $-0.84$ & $-0.06$ & $-0.28$ & 0.94 & 1.36 \\
    \YM & $-0.23$ & 0.34 & 0.44 & $-0.28$ & $-1.67$ & $-0.61$ \\
    \ZN & $-1.24$ & $-1.01$ & $-1.73$ & $-0.52$ & $-0.25$ & $-0.89$ \\
    Equities (\ES\ + \NQ) & 0.71 & $-0.87$ & 0.70 & 0.78 & 1.04 & 1.58 \\
    \quad return & $+1.3\%$ & $-2.5\%$ & $+1.6\%$ & $+1.7\%$ & $+3.1\%$ & $+2.3\%$ \\
    \ES\ + \NQ\ + \YM & 0.46 & $-0.47$ & 0.80 & 0.49 & $-0.01$ & 1.16 \\
    Combined (\ES\ + \NQ\ + \ZN) & 0.71 & $-0.87$ & $-0.23$ & 0.77 & 1.03 & 1.57 \\
    \bottomrule
  \end{tabular}
  \begin{tablenotes}\footnotesize
    \item Notes: 2026 is a partial year (to June).
  \end{tablenotes}
  \end{threeparttable}
\end{table}

\subsection{Regime change for the 10-year}\label{sec:treasury_result}

\ZN\ went through a larger change of regime than the equity indices. Before 2021 policy rates
were near zero, \ZN's moves after releases were small, and a round trip at a quarter tick cost
13--21\% of its typical 30-minute move in every year. The 2022 tightening made rate news the
market's main driver and \ZN's post-release moves larger: the cost ratio fell from 17\% in 2021
to 15\% in 2022 and 10\% in 2023, before rising again to 12--17\% in 2024--2026. Before costs the
book does find something in \ZN: a gross Sharpe ratio of 0.34 and a small positive directional IC
(+0.016). But the regime change did not make the direction tradeable. In 2023, when the larger
moves let \ZN\ through the cost gate for part of the year, it lost, and the combined book scored
$-0.23$ that year against 0.70 for equities. \ZN\ loses in every test year at a quarter tick, its
strategy tracks being always long in the same windows ($-7.4\%$ for both), and it breaks even at
only 0.07 ticks per side. Without a gate, fixed equal-risk weights would have given \ZN\ 60\% of
the combined book at every cost and a Sharpe ratio of $-0.23$; the cost gate is what keeps the
combined book at 0.40. Bigger rate moves did not mean that the past paths forecast the next one:
in \ZN\ the tick is larger than the edge in any regime we observe.

\subsection{Against the surprise strategy}\label{sec:vs_surprise}

The natural competitor to the DTW book is a strategy that takes its direction from the consensus
surprise: buy after a release that beats the median forecast in the direction the market has
rewarded, sell after one that misses. Our companion study runs such books on the same four markets,
with the same releases, an entry one minute after the release and a 30-minute hold, over 2019--2026
and at a harsher cost of two ticks round trip plus \$4.50 in fees. It also runs a simple DTW-only
book at the first window and two books that trade only when the surprise and the DTW forecast agree.
These are not the ladder book of this paper---they use one window, a longer sample and a different
cost---but they show where the direction of the trade is best taken from.


\begin{figure}[!tb]
  \centering
  \begin{minipage}[t]{0.49\linewidth}\centering
    \includegraphics[width=\linewidth]{nb07b_sharpe_by_market}
    \caption{Net Sharpe ratio by market and strategy, 2019--2026, two ticks round trip plus \$4.50.
  ``All'': every release family that passes the impact screen; ``CPI'': CPI releases only;
  ``naive'': trade the sign of the surprise; ``surprise model'': trade a fitted drift forecast;
  ``DTW only'': trade the sign of the DTW forecast at the first window; ``surprise + DTW agree'':
  trade the surprise only when the DTW forecast has the same sign. \ZN\ has no CPI agree trades.}
    \label{fig:vs_surprise}
  \end{minipage}\hfill
  \begin{minipage}[t]{0.49\linewidth}\centering
    \includegraphics[width=\linewidth]{nb07b_portfolio_cumulative}
    \caption{Cumulative net P\&L of the broad surprise book (``All $\cdot$ naive'') by market,
  risk-weighted, and of the two-, three- and four-market portfolios with equal-risk weights fixed on
  2016--2018 (\ES\ 1.00, \NQ\ 0.71, \YM\ 1.02, \ZN\ 2.67), 2019--2026, net of costs.}
    \label{fig:vs_surprise_pnl}
  \end{minipage}
\end{figure}

Figure~\ref{fig:vs_surprise} gives the answer market by market. In all three equity indices the
broad surprise book earns about the same net Sharpe ratio, 0.65 in \ES, 0.71 in \NQ\ and 0.68 in
\YM, and the CPI surprise model earns 0.45--0.76. The DTW-only book is weaker in each: it loses in
\ES\ ($-0.39$), is flat in \NQ\ (0.08) and earns 0.47 in \YM. DTW adds most as a filter on the
surprise. Keeping only the trades on which the two agree raises the Sharpe ratio of the broad book
in \YM\ to 0.94 (0.94 also without 2022), and keeping only the CPI trades on which they agree raises
\ES's CPI book from 0.76 to 0.93, on fewer trades. In \ES\ the broad agree book is lower than the
surprise book alone (0.31), so the filter does not help everywhere. In \ZN\ every book loses after
costs, as in \ref{sec:treasury_result}: the broad surprise book goes from $-0.30$ before costs to
$-2.42$ after them.


Figure~\ref{fig:vs_surprise_pnl} combines the surprise books across markets. The three equity
books rise together, mostly in 2022, and the \ES\ + \NQ\ + \YM\ portfolio earns a Sharpe ratio of
0.78 (37.5\% in total) against 0.72 for \ES\ + \NQ. \ZN\ loses steadily, and because equal-risk
weights give it the largest weight, it turns the four-market portfolio negative ($-0.34$). Both
results match the DTW book: the equity indices carry the edge, adding \YM\ broadens the book, and
\ZN\ should be kept out at these costs. The comparison also locates the DTW book's role. As a
source of direction it is weaker than the surprise in every market; as a second opinion on the
surprise and as a book whose returns are nearly uncorrelated with it, it is useful.

Figure~\ref{fig:vs_surprise_cum} shows how these books earn over time, on all screened releases and
on CPI alone, and adds a fourth book that blends the two signals instead of requiring agreement.
In \ES\ and \NQ\ the surprise-only book is the highest curve on the broad set of releases, and the
DTW-only book drifts down in \ES\ from 2022 and goes flat in \NQ\ after 2023, the same break as in
\ref{sec:2022}. \YM\ is the one market where DTW adds to the surprise: the blend earns a Sharpe
ratio of 1.06 and the agree book 0.94, against 0.68 for the surprise alone, with a smoother path and
without the 2020 drawdown of the surprise book. On CPI the four books are close in every equity
index and most of their P\&L comes from 2022, so the CPI comparison rests on a single year. In
\ZN\ all four curves fall steadily; the surprise book falls most because it trades every release
and pays the cost each time.

\begin{figure}[!tb]
  \centering
  \includegraphics[width=\linewidth]{nb06_cumulative_2x2}
  \caption{Cumulative net P\&L, 2019--2026, two ticks round trip plus \$4.50, of four single-window
  books: surprise only, DTW only, surprise + DTW agree, and a blend of the surprise and DTW signals.
  For each market, left: all screened releases; right: CPI only. \ZN\ has no CPI agree
  trades (flat green line).}
  \label{fig:vs_surprise_cum}
\end{figure}

The surprise books above are built on the broad set of screened releases. Our first walk-forward
study instead trades only the three largest reports---CPI, Core PCE and payrolls---with five
rules: the naive sign of the surprise, a fitted surprise model, and three versions of that model
conditioned on the market's state before the release (volatility, momentum and the like). Every
slope, sign and state model is re-estimated on earlier releases only, and trades pay two ticks plus
\$4.50. On CPI
alone the naive sign is the best or close to the best in each equity index, with Sharpe ratios of
1.03 in \ES, 1.00 in \NQ\ and 0.99 in \YM, and conditioning on the state does not improve on it
by a significant margin. As in the DTW book, almost all of the gain comes in 2022. Adding PCE and payrolls lowers the naive rule to
0.33 in \ES\ and 0.38 in \NQ\ (0.57 in \YM). In \ZN\ the fitted models barely trade, because their forecast rarely exceeds a tick,
and the naive rule loses ($-0.05$ on CPI, $-0.65$ on the three reports). The surprise, in other
words, gives a better direction than the DTW analogues in the equity indices whether it is traded
on one report or many, and neither source survives the \ZN\ tick.


